Weight quantization is the default way to deploy open language models, yet bit widths are usually chosen from perplexity or a single multiple-choice benchmark. We ask which task types lose accuracy first as precision drops, whether the loss is gradual or a cliff, whether smaller models break at higher precision, and whether KL divergence from the full-precision model warns of the cliff in advance. We quantize seven instruction-tuned models (0.5B--8B; \qwen, \llama, \phithree) with weight-only round-to-nearest at 8, 6, 5, 4, 3 and 2 bits (group size 128) and score them on six task types: arithmetic, factual QA, reading comprehension, code, instruction following and creative writing. The cliff lies at or below 4$\to$3 bits for every task, not between INT8 and INT4. Eight and six bits are essentially lossless (3 significant drops in 72 paired tests). At 4 bits the loss is modest but task-dependent: code is the most sensitive task (mean retained score 0.84) and instruction following the least (0.97). Larger models are more robust (Spearman $\rho=-0.53$ between size and critical precision), but not monotonically. Task-specific KL separates broken from healthy configurations with AUC 0.97 at a threshold near 0.21 nats. However, it is a concurrent signal rather than a leading one: one level before failure it exceeds 0.3 nats in only 6 of 42 cases. Aggregate accuracy also hides churn: 7--15\% of correct answers flip at 5 bits, and over 90\% of 4-bit errors are fluent and confident. Deployers of INT4 models should therefore validate each task directly rather than rely on perplexity or a single KL threshold.
